\documentclass[12pt, a4paper]{article}

% --- 常用宏包 ---
\usepackage[UTF8]{ctex}
\usepackage{amsmath, amsthm, amssymb, amsfonts}
\usepackage{geometry}
\usepackage{fancyhdr}
\usepackage{enumerate}
\usepackage{graphicx}
\usepackage{hyperref}
\usepackage{bm}
\usepackage{tcolorbox}

% --- 页面设置 ---
\geometry{left=2.5cm, right=2.5cm, top=2.5cm, bottom=2.5cm}
\pagestyle{fancy}
\fancyhf{}
\rhead{\today}
\lhead{近世代数作业}
\cfoot{\thepage}

% --- 环境定义 ---
\newtcolorbox{problem}[1]{
    colback=gray!5!white,
    colframe=gray!75!black,
    fonttitle=\bfseries,
    title=题目 #1
}

\newenvironment{solution}{
    \begin{proof}[\textbf{解}]
}{
    \end{proof}
}

% --- 个人信息 ---
\title{近世代数作业 07}
\author{李睿阳 (524120910059)}
\date{\today}

\begin{document}
\maketitle


\begin{problem}{1}
    设 $\mathbb{Z}_n$ 是模 $n$ 剩余类环，证明 $I$ 是 $\mathbb{Z}_n$ 的理想 $\iff I$ 是 $\mathbb{Z}_n$ 的加法子群。
\end{problem}
\begin{solution}
    $(\implies)$ 由理想定义，$I$ 是 $\mathbb{Z}_n$ 的加法子群。

    $(\impliedby)$ 设 $I$ 是 $\mathbb{Z}_n$ 的加法子群。对 $\forall r \in \mathbb{Z}_n, a \in I$，
    由于 $r = \underbrace{1 + 1 + \dots + 1}_{r \text{ 个}}$。
    则 $ra = \underbrace{a + a + \dots + a}_{r \text{ 个}}$。
    由加法封闭性，$ra \in I$。由于 $\mathbb{Z}_n$ 是交换环，故 $ar = ra \in I$。
    因此 $I$ 是理想。
\end{solution}


\begin{problem}{2}
    使用题 1 的结论，求 $\mathbb{Z}_6$ 的所有理想。
\end{problem}
\begin{solution}
    由(1) $\mathbb{Z}_6$ 的理想与其加法子群相同，并且由于 $\mathbb{Z}_6$ 为循环群，它的加法子群也都是循环群。
    它的理想为：
    \begin{itemize}
        \item $\langle 0 \rangle = \{0\}$
        \item $\langle 1 \rangle = \{0, 1, 2, 3, 4, 5\}$
        \item $\langle 2 \rangle = \{0, 2, 4\}$
        \item $\langle 3 \rangle = \{0, 3\}$
    \end{itemize}
\end{solution}


\begin{problem}{3}
    设 $\langle a \rangle, \langle b \rangle$ 是整数环 $\mathbb{Z}$ 的两个理想，证明：
    (1) $\langle a \rangle \cap \langle b \rangle = \langle [a, b] \rangle$
    (2) $\langle a \rangle + \langle b \rangle = \langle (a, b) \rangle$
\end{problem}
\begin{solution}
    (1) $k \in \langle a \rangle \cap \langle b \rangle \iff a|k$ 且 $b|k \iff [a, b] | k \iff k \in \langle [a, b] \rangle$。
    故 $\langle a \rangle \cap \langle b \rangle = \langle [a, b] \rangle$。

    (2) 因为 $(a, b) | a$ 且 $(a, b) | b$，故 $a, b \in \langle (a, b) \rangle$，从而 $\langle a \rangle + \langle b \rangle \subseteq \langle (a, b) \rangle$。
    反之，由裴蜀定理，存在 $x, y \in \mathbb{Z}$ 使得 $(a, b) = xa + yb$。
    故 $(a, b) \in \langle a \rangle + \langle b \rangle$，从而 $\langle (a, b) \rangle \subseteq \langle a \rangle + \langle b \rangle$。
    故 $\langle a \rangle + \langle b \rangle = \langle (a, b) \rangle$。
\end{solution}
    

\begin{problem}{4}
    在环 $\mathbb{Z}[\sqrt{-5}]$ 中，设 $I = \langle 3, 1 + \sqrt{-5} \rangle$，证明对任意的 $a + b\sqrt{-5} \in I$，有 $3 | a - b$。
\end{problem}
\begin{solution}
    对任意 $a + b\sqrt{-5} \in I$，存在 $x + y\sqrt{-5}, u + v\sqrt{-5} \in \mathbb{Z}[\sqrt{-5}]$，使得
    \begin{align*}
        a + b\sqrt{-5} &= 3(x + y\sqrt{-5}) + (1 + \sqrt{-5})(u + v\sqrt{-5}) \\
        &= (3x + u - 5v) + (3y + u + v)\sqrt{-5}
    \end{align*}
    由此可得：
    $a = 3x + u - 5v, \quad b = 3y + u + v$
    则 $a - b = (3x + u - 5v) - (3y + u + v) = 3x - 3y - 6v = 3(x - y - 2v)$。
    所以 $3 | a - b$。
\end{solution}


\end{document}
